\section{The full kernel with a fixed parent angle}\label{dt:sec-2}

The fixed Newton polynomials are

\begin{equation}\label{dt:eq-N1}
\begin{gathered}
 d=N-1,\qquad c_j=2\cos((2j+1)\pi/N),\qquad
 S_k=\prod_{r=1}^k2\sin(\pi r/N),\\
 e_k(z)=S_k^{-2}\prod_{j<k}(c_j-z),\qquad 0\le k<N.
 \end{gathered}
\end{equation}

Let \(\mathcal T_N(2\cos\theta)=2\cos(N\theta)\), and put
\(P_\sigma(z)=\mathcal T_N(z)+2\cos(2\pi\sigma)\). The normalized
polynomial trace is

\begin{equation}\label{dt:eq-N2}
\phi_\sigma(F)=(-1)^dN^2[z^d]\operatorname{rem}_{P_\sigma}F,
 \qquad \phi_\sigma(e_i e_j)=\mathbf1_{i+j=d}.
\end{equation}

The remainder has degree less than \(N\). This is the full polynomial
functional, including its derivative values when roots coalesce.
\cite[Theorem~7.1]{Meridian}
retains the paired factors

\begin{equation}\label{dt:eq-N3}
W_k(s)=\prod_{r=1}^k
     \left(1-\frac{\sin^2(\pi s/N)}{\sin^2(\pi r/N)}\right).
\end{equation}

After lowering the pairing \eqref{dt:eq-N2}, a physical turn gives the equal-index
weight \(\gamma_e(s)H_N(s)/W_k(s)\). The factor \(\gamma_e\) is the
original scalar prescribed by that edge's orientation type. It stays in
the exact scalar ledger.

For the edge action write

\begin{equation}\label{dt:eq-N4}
\begin{gathered}
 J(x)=\int_0^x\log(2\sin\pi v)\,dv,\qquad
 g(x)=J(2x)-4J(x),\qquad G_0=g(1/4),\\
 b(x,t)=\int_0^x\log(2\cos(2\pi v)+2\cos t)\,dv-2J(x),\\
 \psi(t,u)=\operatorname{stat}_x[b(x,t)+b(x,u)],
                       \qquad\psi(0,0)=2G_0.
 \end{gathered}
\end{equation}

The logarithms and stationary branch are continued from
\(x=1/4,t=u=0\). The canonical factorial interpolation is \cite[Lemma~2.1]{Factorial}.
Its midpoint application and the finite complex edge contour are proved
in \cite[Section~3]{Summed}. The full polynomial and repeated-root
bounds are \cite[Lemma~2.1]{Summed}. These statements precede the
summed hypothesis in their proofs and hold independently of it.
Use the exact full kernel

\begin{equation}\label{dt:eq-13}
K_N(t,u;s)=\sum_{k=0}^{N-1}
       e_k(-2\cos t)e_k(-2\cos u)\frac{H_N(s)}{W_k(s)}.
\end{equation}

The definitions of \(e_k,W_k\) and the fixed anti-diagonal trace metric are
unchanged. It is a polynomial of degree at most \(N-1\) in each of
\(-2\cos t\) and \(-2\cos u\), hence has trigonometric degree at most
\(N-1\) in either angle.
The paired factors and their inverses are bounded for fixed \(|s|\le R<1\).
Every inequality below permits a polynomial factor in N.

The full-kernel real estimate is

\begin{equation}\label{dt:eq-14}
|K_N(t,u;s)|\le C_RN^{A_R}
   e^{N\{2G_0-k_0[(1-\cos t)+(1-\cos u)]\}},
                 \qquad t,u\in[-\pi,\pi].
\end{equation}

To verify this, split the sum at \(I=[1/5,3/10]\).
On \(I\), \cite[equation~(4.1)]{Summed} gives angular loss
\(\sin^2(t/2)/160\) at each end. The positive zero-angle sum
has rate \(2G_0\). Outside \(I\), the full polynomial bound
has loss \(\eta_K=\pi/200\), including high-index derivatives.
Since \(4k_0<\eta_K\) and \(2k_0<1/160\), both parts obey
\eqref{dt:eq-14}. The bounded paired insertion changes only its
polynomial prefactor.

The conditional estimate needed here is

\begin{equation}\label{dt:eq-15}
|K_N(t+iy,u;s)|\le C_RN^{A_R}
 e^{N\{\Re\psi(0,u)-k_0(1-\cos t)+16y^2\}},
 \quad u\in\mathcal R,\quad t\in[-\pi,\pi],\quad |y|\le b_0.
\end{equation}

Thus the parent compact is a long, thin rectangle, rather than the old
disk of radius \(2^{-12}\). The constant term is the actual conditional phase
\(\psi(0,u)\), not its upper bound \(2G_0+O(|u|^2)\).

Here is a proof from the original finite kernel. Split the indices into
\(I=[1/5,3/10]\) and its complement. The complementary full polynomial, by \cite[Lemma~2.1 and equation~(2.7)]{Summed}, has the single-edge loss \(\eta_K=\pi/200\), including high
indices and their complete derivatives. Its ellipse bound is

\begin{equation}\label{dt:eq-16}
C_RN^{A_R}
       e^{N\{2G_0-\eta_K+|y|+|\Im u|\}}.
\end{equation}

The edge Hessian in \cite[equation~(4.3)]{Summed} gives, on the radius 1/4 bidisk,

\begin{equation}\label{dt:eq-17}
|\psi_{tt}+a_0|\le a_0/50,
 \quad |\psi_{uu}+a_0|\le a_0/50,
 \quad |\psi_{tu}|\le a_0/50,
 \quad \psi_t(0,u)=\psi_u(t,0)=0.
\end{equation}

In particular \(|\psi(0,u)-2G_0|\le |u|^2/10\). Hence \eqref{dt:eq-16} is smaller
than \eqref{dt:eq-15}: the sum \(b_0+b_0+(R_0^2+b_0^2)/10+2k_0\) is less
than \(\eta_K/2\). All these constants are independent of the tree.

For the retained indices write \(u=a+ib\). Their factors at \(a\) are positive.
If \(r_j=\cos(2\pi(j+1/2)/N)\), then \(r_j>-1/3\) and
\(r_j+\cos a>253/384\). Along \(|a|\le1/8,\ |b|\le2^{-11}\), the denominator
\(r_j+\cos(a+ib)\) has modulus greater than \(3/5\), and

\begin{equation}\label{dt:eq-18}
\left|\frac{d^2}{du^2}\log(r_j+\cos u)\right|
 =\left|\frac{-r_j\cos u-1}{(r_j+\cos u)^2}\right|<9.
\end{equation}

The first derivative at real a is real. Integrating vertically therefore
gives the actual finite-product estimate

\begin{equation}\label{dt:eq-19}
|e_k(-2\cos(a+ib))|
       \le e_k(-2\cos a)e^{5Nb^2},\qquad k/N\in I.
\end{equation}

At the other end retain the original product bound

\begin{equation}\label{dt:eq-20}
|e_k(-2\cos(t+iy))|
   \le e^2e_k(-2)e^{N[-\alpha\sin^2(t/2)+11y^2]},
                       \qquad \alpha=1/160>1/250.
\end{equation}

The positive retained sum at \((0,a)\) is bounded by a polynomial times
\(\exp(N\psi(0,a))\), uniformly for real \(|a|\le1/8\). This is the positive one-edge Laplace estimate in \cite[Section~3]{Summed}: its unique critical wire is

\begin{equation}\label{dt:eq-21}
\tan^2(\pi x)=\frac1{1+\tan^2(a/2)},
\end{equation}

which stays in the interior of \(I\). Its second derivative stays negative
on this fixed real compact. Equations \eqref{dt:eq-19}--\eqref{dt:eq-20}, and the second bound
in \eqref{dt:eq-17}, consequently give

\begin{equation}\label{dt:eq-22}
|K_{N,I}(t+iy,u;s)|\le C_RN^{A_R}
 e^{N\{\Re\psi(0,u)-\alpha\sin^2(t/2)+11y^2+6b^2\}}.
\end{equation}

For \(|t|\ge1/8\),

\begin{equation}\label{dt:eq-23}
6b_0^2<\frac{\alpha}{2}\sin^2(1/16),\qquad
                         2k_0<\alpha/2.
\end{equation}

For example \(\sin(1/16)>1/(8\pi)\), \(\pi^2<10\) and \(\alpha>1/250\) already
prove the first inequality. This proves \eqref{dt:eq-15} in this range.

For \(|t|\le1/8\), \cite[equation~(3.10)]{Summed} gives the exact
nonzero edge asymptotic on the 1/4 bidisk. Both \(u\) and \(t+iy\) lie strictly
inside that disk. Integrating the first bound of \eqref{dt:eq-17} gives \eqref{dt:eq-15} there.
The original lattice currents and \eqref{dt:eq-16} are exponentially smaller.
In particular the full kernel, uniformly on this smaller rectangle
product, retains

\begin{equation}\label{dt:eq-24}
K_N(t,u;s)=N^{-3/2}a_e(t,u)e^{N\psi(t,u)}
                         (1+O_R(N^{-1})),
\end{equation}

with the original nonzero holomorphic leading amplitude and midpoint
normalization. The outside physical edge gamma is not included in \(a_e\).
